% group: scales-arpeggios
\begin{patternbox}{Diatonic seventh arpeggios (thirds through the scale)}{the major key: \ch{I}{\Delta7} \ch{ii}{-7} \ch{iii}{-7} \ch{IV}{\Delta7} \ch{V}{7} \ch{vi}{-7} \ch{vii}{\text{\o}7}}
The arpeggio \mel{1\;3\;5\;7} built on every degree of the major scale, in thirds,
one bar per chord:
\mel{1\;3\;5\;7} $|$ \mel{2\;4\;6\;\hi{1}} $|$ \mel{3\;5\;7\;\hi{2}} $|$ \mel{4\;6\;\hi{1}\;\hi{3}} $|$
\mel{5\;7\;\hi{2}\;\hi{4}} $|$ \mel{6\;\hi{1}\;\hi{3}\;\hi{5}} $|$ \mel{7\;\hi{2}\;\hi{4}\;\hi{6}} $|$ \mel{\hi{1}}.
Guitar: two octaves of arpeggios from the low tonic (low E string, 3rd fret in G) up
to the high tonic. Stay in position (frets 2--5) as long as possible, then move up the
neck on the top three strings.
\patterntab{1,@6 3,@5 5,@5 7,@4 | 2,@6 4,@5 6,@4 1@4 | 3,@5 5,@5 7,@4 2@3 | 4,@5 6,@4 1@4 3@3}
\patterntab{5,@5 7,@4 2@3 4@3 | 6,@4 1@4 3@3 5@2 | 7,@4 2@3 4@3 6@2 | 1@4 3@3 5@2 7@1}
\patterntab{2@3 4@3 6@2 1'@1 | 3@3 5@2 7@1 2'@1 | 4@3 6@2 1'@2 3'@1 | 5@3 7@2 2'@1 4'@1}
\patterntab{6@3 1'@2 3'@1 5'@1 | 7@3 2'@2 4'@2 6'@1 | 1''@1}
\patternwhereline
\end{patternbox}
